% GENERATED FILE. Edit canonical lessons, then run npm run generate:books.
% canonical-source-hash: fnv1a64:739b08cc257f142c
% canonical-lessons: MR-C203-nakki, MR-C203-kharokhar, MR-C203-javaljaval, MR-C203-sumare, MR-C203-phakta

\chapter{Certainly, Really, Almost, About, Only}
\label{ch:mr-certainly-really-almost-about-only}
\hlchaptermodality{203}

\begin{chapteropening}
\textbf{By the end of this chapter:} \emph{I can say certainly, really, almost, about and only.}

Complete the last lesson of chapter 203: I can say certainly, really, almost, about and only.
\end{chapteropening}

\section[\texorpdfstring{\mr{नक्की} (nakkī)}{नक्की (nakkī)}]{\mr{नक्की} (nakkī) — certainly, for sure}

\label{lesson:MR-C203-nakki}



\begin{quote}
Before the new one: say the Marathi for a suburb, then the Marathi for up.
\end{quote}

\subsection*{You'll want to know: \mr{नक्की}}
\textbf{\mr{नक्की}} — \emph{nakkī} — \textquotedblleft{}certainly, for sure\textquotedblright{}.

\begin{cousinweb}
The origin of \textbf{\mr{नक्की}} (\emph{nakkī}) is not settled. Neither CDIAL nor DEDR, searched by form and by meaning, has an entry for it.
\end{cousinweb}

\begin{grammarlens}[title={What you've built}]
One more word for this chapter.
\end{grammarlens}

\subsection*{Your turn}
\begin{itemize}
\raggedright
  \item \emph{Say it:} \emph{nakkī}
  \item \emph{Say it:} \emph{nakkī}, once more
  \item \emph{Say it:} say \emph{nakkī}
  \item \emph{Recall:} say the Marathi for a suburb, then the Marathi for up, then say \emph{nakkī} again
\end{itemize}

\begin{tcolorbox}[breakable,colback=teal!4,colframe=teal!35!black,title={Before you move on}]
What does \mr{नक्की} mean? (\textquotedblleft{}Certainly, for sure\textquotedblright{}.) Say it once more.
\end{tcolorbox}
\section[\texorpdfstring{\mr{खरोखर} (kharokhar)}{खरोखर (kharokhar)}]{\mr{खरोखर} (kharokhar) — really}

\label{lesson:MR-C203-kharokhar}



\begin{quote}
Before the new one: say the Marathi for up, then the Marathi for certainly.
\end{quote}

\subsection*{You'll want to know: \mr{खरोखर}}
\textbf{\mr{खरोखर}} — \emph{kharokhar} — \textquotedblleft{}really\textquotedblright{}.

\begin{cousinweb}
\textbf{\mr{खरोखर}} (\emph{kharokhar}) is \emph{kharā}, \textquotedblleft{}true\textquotedblright{}, doubled with a linking \textbf{-\mr{ो}-}. CDIAL lists Marathi \emph{kharā}, \textquotedblleft{}pure, good, firm\textquotedblright{}, under Sanskrit \emph{khara}, \textquotedblleft{}hard, sharp\textquotedblright{} (CDIAL 3819).
\end{cousinweb}

\begin{grammarlens}[title={What you've built}]
One more word for this chapter.
\end{grammarlens}

\subsection*{Your turn}
\begin{itemize}
\raggedright
  \item \emph{Say it:} \emph{kharokhar}
  \item \emph{Say it:} \emph{kharokhar}, once more
  \item \emph{Say it:} say \emph{kharokhar}
  \item \emph{Recall:} say the Marathi for up, then the Marathi for certainly, then say \emph{kharokhar} again
\end{itemize}

\begin{tcolorbox}[breakable,colback=teal!4,colframe=teal!35!black,title={Before you move on}]
What does \mr{खरोखर} mean? (\textquotedblleft{}Really\textquotedblright{}.) Say it once more.
\end{tcolorbox}
\section[\texorpdfstring{\mr{जवळजवळ} (javaḷjavaḷ)}{जवळजवळ (javaḷjavaḷ)}]{\mr{जवळजवळ} (javaḷjavaḷ) — almost}

\label{lesson:MR-C203-javaljaval}



\begin{quote}
Before the new one: say the Marathi for certainly, then the Marathi for really.
\end{quote}

\subsection*{You'll want to know: \mr{जवळजवळ}}
\textbf{\mr{जवळजवळ}} — \emph{javaḷjavaḷ} — \textquotedblleft{}almost\textquotedblright{}.

\begin{cousinweb}
\textbf{\mr{जवळजवळ}} (\emph{javaḷjavaḷ}) is \textbf{\mr{जवळ}}, \textquotedblleft{}near\textquotedblright{}, said twice: \textquotedblleft{}near-near\textquotedblright{}. The origin of \textbf{\mr{जवळ}} is not settled; CDIAL has no entry for it.
\end{cousinweb}

\begin{grammarlens}[title={What you've built}]
One more word for this chapter.
\end{grammarlens}

\subsection*{Your turn}
\begin{itemize}
\raggedright
  \item \emph{Say it:} \emph{javaḷjavaḷ}
  \item \emph{Say it:} \emph{javaḷjavaḷ}, once more
  \item \emph{Say it:} say \emph{javaḷjavaḷ}
  \item \emph{Recall:} say the Marathi for certainly, then the Marathi for really, then say \emph{javaḷjavaḷ} again
\end{itemize}

\begin{tcolorbox}[breakable,colback=teal!4,colframe=teal!35!black,title={Before you move on}]
What does \mr{जवळजवळ} mean? (\textquotedblleft{}Almost\textquotedblright{}.) Say it once more.
\end{tcolorbox}
\section[\texorpdfstring{\mr{सुमारे} (sumāre)}{सुमारे (sumāre)}]{\mr{सुमारे} (sumāre) — about, approximately}

\label{lesson:MR-C203-sumare}



\begin{quote}
Before the new one: say the Marathi for really, then the Marathi for almost.
\end{quote}

\subsection*{You'll want to know: \mr{सुमारे}}
\textbf{\mr{सुमारे}} — \emph{sumāre} — \textquotedblleft{}about, approximately\textquotedblright{}.

\begin{cousinweb}
\textbf{\mr{सुमारे}} (\emph{sumāre}) is often traced to Persian \emph{shumār}, \textquotedblleft{}count, number\textquotedblright{}, but no source here settles it. Neither CDIAL nor DEDR has an entry for it.
\end{cousinweb}

\begin{grammarlens}[title={What you've built}]
One more word for this chapter.
\end{grammarlens}

\subsection*{Your turn}
\begin{itemize}
\raggedright
  \item \emph{Say it:} \emph{sumāre}
  \item \emph{Say it:} \emph{sumāre}, once more
  \item \emph{Say it:} say \emph{sumāre}
  \item \emph{Recall:} say the Marathi for really, then the Marathi for almost, then say \emph{sumāre} again
\end{itemize}

\begin{tcolorbox}[breakable,colback=teal!4,colframe=teal!35!black,title={Before you move on}]
What does \mr{सुमारे} mean? (\textquotedblleft{}About, approximately\textquotedblright{}.) Say it once more.
\end{tcolorbox}
\section[\texorpdfstring{\mr{फक्त} (phakta)}{फक्त (phakta)}]{\mr{फक्त} (phakta) — only}

\label{lesson:MR-C203-phakta}



\begin{quote}
Before the new one: say the Marathi for almost, then the Marathi for about.
\end{quote}

\subsection*{You'll want to know: \mr{फक्त}}
\textbf{\mr{फक्त}} — \emph{phakta} — \textquotedblleft{}only\textquotedblright{}.

\begin{cousinweb}
\textbf{\mr{फक्त}} (\emph{phakta}) is a loan, usually traced through Persian to Arabic \emph{faqaṭ}, \textquotedblleft{}only\textquotedblright{}, with \textbf{\mr{फ}} for the \emph{f}.
\end{cousinweb}

\begin{grammarlens}[title={What you've built}]
One more word for this chapter.
\end{grammarlens}

\subsection*{Your turn}
\begin{itemize}
\raggedright
  \item \emph{Say it:} \emph{phakta}
  \item \emph{Say it:} \emph{phakta}, once more
  \item \emph{Say it:} say \emph{phakta}
  \item \emph{Recall:} say the Marathi for almost, then the Marathi for about, then say \emph{phakta} again
\end{itemize}

\begin{tcolorbox}[breakable,colback=teal!4,colframe=teal!35!black,title={Before you move on}]
What does \mr{फक्त} mean? (\textquotedblleft{}Only\textquotedblright{}.) Say it once more.
\end{tcolorbox}
